\documentclass[12pt]{article}
%\usepackage{alltt}
%\usepackage{helvet}
%\usepackage[sfdefault]{roboto}
\usepackage{amsmath}
\usepackage[utf8]{inputenc}
\usepackage[dvips]{graphicx}
%\usepackage{a4wide}
\usepackage{epsfig}
\usepackage{fancybox}
\usepackage{verbatim}
\usepackage{array}
\usepackage{latexsym}
\usepackage{alltt}
\usepackage{url}
\usepackage{color}   % added for colored text
%\usepackage{fullpage}
\usepackage{hyperref}
\usepackage{listings}
\usepackage{color}
\usepackage{calc}
\usepackage{enumitem}
\usepackage[hmargin=3cm,vmargin=5.0cm]{geometry}
%\topmargin=0cm
\topmargin=-1.8cm \addtolength{\textheight}{4.5cm}
\addtolength{\textwidth}{1.0cm}
%\setlength{\leftmargin}{-5cm}
\setlength{\oddsidemargin}{0.0cm}
\setlength{\evensidemargin}{0.0cm}

%\renewcommand{\familydefault}{\sfdefault}

\newcommand{\HRule}{\rule{\linewidth}{1mm}}
\newcommand{\kutu}[2]{\framebox[#1mm]{\rule[-2mm]{0mm}{#2mm}}}
\newcommand{\gap}{ \\[1mm] }

\newcommand{\Q}{\raisebox{1.7pt}{$\scriptstyle\bigcirc$}}

\definecolor{amaranth}{rgb}{0.9, 0.17, 0.31}
\definecolor{gray}{rgb}{0.4,0.4,0.4}
\definecolor{darkblue}{rgb}{0.0,0.0,0.6}
\definecolor{cyan}{rgb}{0.0,0.6,0.6}
\definecolor{red}{rgb}{0.6,0,0}
\definecolor{dkgreen}{rgb}{0,0.6,0}
\definecolor{mauve}{rgb}{0.58,0,0.82}
\definecolor{lightblue}{rgb}{0.0,0.0,0.9}
\definecolor{darkred}{rgb}{0.6,0.0,0.0}

\lstset{
    %backgroundcolor=\color{lbcolor},
    tabsize=4,
    basicstyle=\fontsize{10}{10.3}\selectfont\sffamily,
    numberstyle=\footnotesize,
    aboveskip={0.0\baselineskip},
    belowskip={0.0\baselineskip},
    columns=fullflexible,
    breaklines=true,
    prebreak=\raisebox{0ex}[0ex][0ex]{\ensuremath{\hookleftarrow}},
    frame=single,
    showtabs=false,
    showspaces=false,
    showstringspaces=false,
    identifierstyle=\color{amaranth},
    keywordstyle=\color{rgb}{0,0,1},
    commentstyle=\color[rgb]{0.133,0.545,0.133},
    stringstyle=\color{amaranth},
}

\lstdefinelanguage{XML}
{
  morestring=[s][\color{red}]{"}{"},
  morestring=[s][\color{black}]{>}{<},
  morecomment=[s]{<?}{?>},
  morecomment=[s][\color{dkgreen}]{<!--}{-->},
  %stringstyle=\color{black},
  identifierstyle=\color{black},
  keywordstyle=\color{darkblue},
  commentstyle=\color[rgb]{0.133,0.545,0.133},
  stringstyle=\color{black},
  morekeywords={Scene,BackgroundColor,ShadowRayEpsilon,
  MaxRecursionDepth,Cameras,Camera,Position,Gaze,Up,
  NearPlane,NearDistance,ImageResolution,ImageName,
  Material,Materials,VertexData,Mesh,Triangle,Sphere,
  Faces,Lights,AmbientLight,PointLight,Objects,Indices,
  AmbientReflectance,DiffuseReflectance,SpecularReflectance,PhongExponent,MirrorReflectance,Intensity,Center,Radius,xmlns,version}% list your attributes here
}

\begin{document}

\noindent \HRule \\[3mm]
\small
\begin{tabular}[b]{lp{1.2cm}r}
\href{https://www.metu.edu.tr/}{\epsfig{file=metulogo.eps,width=5mm}} Middle East Technical
University &  &
\href{https://ceng.metu.edu.tr/information}{\epsfig{file=bmblogo.eps,width=5mm}} Department of Computer Engineering \\
\end{tabular} \\
\begin{center}

                 \LARGE \textbf{CENG 795} \\[4mm]
                 \Large Advanced Ray Tracing \\[4mm]
                \normalsize Fall '2022-2023 \\
                    \normalsize Assignment 2 - Speed and Geometry \\
                    \normalsize (v.1.0)
\end{center}
\HRule

\begin{center}
Due date: November 13, 2022, Sunday, 23:59
\end{center}


\section{Objectives}
The focus of this assignment is to add support for acceleration
structures, geometric transformations, and instancing to your ray
tracers. These features will help you to create and render much more
complex scenes.

\vspace{0.5cm} \noindent \textbf{Keywords:} \emph{transformations,
    instancing, bounding volume hierarchy (BVH), kd-tree}

\section{Specifications}
\begin{enumerate}

\item \textbf{}You should name your executable as ``raytracer''.

\item \textbf{}Your executable will take an XML scene file as argument
(e.g. ``scene.xml''). A parser will be given to you, so that you do not
have to worry about parsing the file yourself. The format of the file
will be explained in Section \ref{sec:sceneFile}. You should be able to
run your executable via command  ``./raytracer scene.xml''.

\item \textbf{}You will save the resulting images in the PNG format. You
can use a library of your choice for saving PNG images.

\item \textbf{}The scene file may contain multiple camera
configurations. You should render as many images as the number of
cameras. The output filenames for each camera is also specified in the
XML file.

\item \textbf{}There is no time limit for rendering the input scenes.
However, if you implement an acceleration structure such as kd-tree or
BVH correctly, it should take only a few seconds to render even the most
complex scenes.

\item \textbf{}You should use Blinn-Phong shading model for the specular
shading computations.  Mirrors and dielectrics should obey Fresnel
reflection rules.  Dielectrics are assumed to be isotropic and should
obey Beer's law.

\item \textbf{}You will implement two types of light sources: point and
ambient. There may be multiple point light sources and a single ambient
light. The values of these lights will be given as (R, G, B) color
triplets that are \underline{not} restricted to $[0, 255]$ range (however,
        they cannot be negative as negative light does not make sense).
Any pixel color value that is calculated by shading computations and is
greater than 255 must be clamped to 255 and rounded to the nearest
integer before writing it to the output PNG file. This step will be
replaced by the application of a tone mapping operator in our later
homeworks.

\item \textbf{}Point lights will be defined by their intensity (power
        per unit solid angle). The irradiance due to such a light source
falls off as inversely proportional to the squared distance from the
light source. To simulate this effect, you must compute the irradiance
at a distance of $d$ from a point light as:
%
\begin{center}
$E(d) = \frac{I}{d^2}$,
\end{center}
%
where $I$ is the original light intensity (a triplet of RGB values given
in the XML file) and $E(d)$ is the irradiance at distance $d$
from the light source.

\item \textbf{Back-face culling} is a method used to accelerate the ray
- scene intersections by not computing intersections with triangles
whose normals are pointing away from the camera. Its implementation is
simple and done by calculating the dot product of the ray direction with
the normal vector of the triangle. If the sign of the result is
positive, then that triangle is ignored. Note that shadow rays should
not use back-face culling. In this homework, back-face culling
implementation is optional. Experiment with enabling and disabling it in
your ray tracers and report your time measurements in your blog post.

\item \textbf{Degenerate triangles} are those triangles whose at least
two vertices coincide.  The input files given to you should be free of
such cases, but models downloaded from the Internet occasionally have
this problem. You can put a check in your ray tracer either to detect or
ignore such triangles (using such triangles usually produce NaN values).

\end{enumerate}

\section{Scene File}
\label{sec:sceneFile}

\noindent Please see HW1 for a detailed description of the scene file.
In this section, only the elements introduced in this homework are
explained.

\begin{itemize}

\item \textbf{Transformations:} All transformations are defined in this
element group. Each of the transformations defined below has an id for
referring to them later.

\begin{description}[leftmargin=!,labelwidth=\widthof{\bfseries Translation},labelindent=\parindent]

\item[Translation]  $\Delta x$, $\Delta y$, and $\Delta z$
values corresponding to translation amounts in each axis.

\item[Scaling]  $S_x$, $S_y$, and $S_z$ scaling factors along
each axis.

\item[Rotation]  $\theta$, $x$, $y$, and $z$ values that define the angle of
rotation in degrees and the rotation axis. Positive angles correspond to
counter-clockwise rotations.

\end{description}

%
\item \textbf{AmbientLight:}  is defined by just an X, Y, Z radiance
triplet. This is the amount of light received by each object even when
the object is in shadow. Color channel order of this triplet is RGB.

\item \textbf{PointLight:} is defined by a position and an intensity,
    which are all floating point numbers. The color channel order of
    intensity is RGB.

\item \textbf{Mesh:}  The mesh definition now has an extra element
called \textbf{\textless Transformations\textgreater} that defines which
transformations will apply to this mesh and in which order. For example,
if the value of this element is \textbf{r1 t1 s2 r2}, it means first
apply rotation transformation with id equal to 1, then translation with
id equal to 1, then scaling with id equal to 2 and then rotation with
id equal to 2. If we represent it as a matrix multiplication, the
composite transformation should look like this:

\begin{align}
M = R_2 S_2 T_1 R_1,
\end{align}

and the transformed coordinates of a vertex $v$ will be equal to:

\begin{align}
v' = M v
\end{align}

\item \textbf{Triangle:} The same transformations element can be found
under the triangle element as well.

\item \textbf{Sphere:} The same transformations element can be found
under the sphere element as well.

\item \textbf{MeshInstance:} This is a new object that is introduced by
this homework. A mesh instance is a copy of its base mesh with a
potentially different material and transformation. A mesh instance can
be instanced from another mesh instance as well and because of that ids
of meshes and mesh instances must not be the same. A mesh instance does
not have \textbf{\textless Faces\textgreater} element as it uses the same faces
with its mother mesh. An example usage of this element is given below:

\begin{verbatim}
<MeshInstance id="3" baseMeshId="2" resetTransform="true">
    <Material>3</Material>
    <Transformations>t2 r1</Transformations>
</MeshInstance>
\end{verbatim}

This defines a mesh instance whose mother mesh (or mesh instance) is the
object with id equal to 2. This instance is using material 2, which may
be different from the base object's material. The \textbf{resetTransform} 
attribute, whose default value is \textbf{false}, indicates
whether the defined transformation should apply on top of the
transformations that are already defined for the base object. If its
value is true, the defined transformation is the only one that should
apply on the initial geometry. Otherwise, the new transformation should
be applied on top of what is already defined for the base mesh. This can
be represented as follows:

\begin{align}
M_{instance} = 
\begin{cases}
R_1 T_2 & \textrm{if resetTransform is true}\\
R_1 T_2 M_{base} & \textrm{if resetTransform is false (default)}
\end{cases},
\end{align}

where $M_{base}$ is the composite transformation of the base object.

\end{itemize}

\section{Hints \& Tips}
In addition to those for HW1, the following tips may be useful for this
homework.

\begin{enumerate}
\item \textbf{} There are no special elements that are defined
for acceleration structures. You are expected to implement a structure
of your choice, with BVH being the recommended one.

\item \textbf{} As our goal is to speed up our ray tracers, adding
multi-threading support could be quite beneficial at this point. There
are different ways of implementing this. You can divide the
image region into sub-regions whose count is equal to the number of
threads. Each thread renders its own region. However, if some threads
finish their jobs earlier than others, this may be sub-optimal.
Alternatively, you can divide the image into a larger number of regions
and assign threads to each region. Whenever a thread finishes processing
its own region, it can get assigned to a new one. This would be
somewhat more complex to implement but would result in a more efficient
ray tracer.

\end{enumerate}

\section{Bonus}

I will be more than happy to give bonus points to students who
make important contributions such as new scenes, importers/exporters
between our XML format and other standard file formats. Note that a
Blender
exporter\footnote{\url{https://saksagan.ceng.metu.edu.tr/courses/ceng477/student/ceng477exporter.py}},
which exports Blender data to our XML format, was written by one of our
previous students. You can use this for designing a scene in Blender and
exporting it to our file format.

\section{Regulations}

\begin{enumerate}

\item \textbf{Programming Language:} C/C++ is the recommended language.
However, other languages can be used if so desired. In the past, some
some students used Rust or even Haskell for implementing their ray
tracers.

\item \textbf{Changing the Sample Codes:} You are free to modify any
sample code provided with this homework.

\item \textbf{Additional Libraries:} If you are planning to use any
library other than \textit{(i)} the standard library of the language,
\textit{(ii)} pthread, \textit{(iii)} the XML parser, and the PNG
libraries please first ask about
it on ODTUClass and get a confirmation. Common sense rules apply: if a
library implements a ray tracing concept that you should be
implementing yourself, do not use it!

\item \textbf{Submission:} Submission will be done via ODTUClass. 
To submit, Create a
``\textbf{tar.gz}''  file  named  ``raytracer.tar.gz'' that
contains all your source code files and a Makefile. The
executable should  be  named as  ``raytracer'' and  should  be
able  to  be  run  using  the following commands (scene.xml
        will be provided by us during grading):\\

\indent \textbf{tar -xf raytracer.tar.gz}\\
\indent \textbf{make}\\
\indent \textbf{./raytracer scene.xml}\\

\noindent Any error in these steps will cause point penalty during
grading.

\item \textbf{Late Submission:} You can submit your codes up to 3 days
late. Each late day will cause a 10 point penalty.

\item \textbf{Cheating:} \textbf{We have zero tolerance policy
for cheating}.  People involved in cheating will be
punished according to the university regulations and will
get 0 from the homework. You can discuss algorithmic choices,
but sharing code between groups or using third party code
is strictly forbidden. By the nature of this class, many past students
make their ray tracers publicly available. You must refrain from using
them at all costs.

\item \textbf{Forum:} Check the ODTUClass forum regularly for
updates/discussions.

\item \textbf{Evaluation:} The basis of evaluation is your blog posts.
Please try to create interesting and informative blog posts about your
ray tracing adventures. You can check out various past blogs for
inspiration. However, also expect your codes to be compiled and tested
on some examples for verification purposes. So the images that you share
in your blog post must directly correspond to your ray tracer outputs.

\end{enumerate}

\end{document}
